% Nouveau EGA V (cote 25): archive cote 25, site page 178 (PDF page 179).
% typescript. Draft compared: none.
% AI-transcribed by Claude Opus 5.5 (claude-opus-5-5, adaptive thinking, effort high), Message Batches API, batch ega25b (msgbatch_01A7n5YuXubUPNK1zsepi8RG), 2026-11-13; not yet checked by eye.
\scanpage{25/178}
\uline{Corollaire} 16.10. Soient k un corps, E un vectoriel de dimensionfinie
sur k, $Q \in \mathrm{Sym}^2(E)$. Conditions équivalentes:

(i) Q est lisse.

(ii) $\overset{\hand{$\sim$}}{d}(Q) \neq 0$.

(iii) Il existe une extension \xout{finie} k' de k telle que $Q\otimes k'$ soit
isomorphe à la forme standart.

(iii bis) Comme (iii), avec k' une extension finie de k.

C'est une conséquence triviale de 16.9. On notera que lorsque k est
algébriquement clos, alors (iii) et (iii bis) se remplacent par la
condition plus frappante: Q est isomorphe à la forme standart.

Démonstration de 16.9. On \xout{a évidemment} peut évidemment
supposer E de rang constant n. On aura évidemment
\[ \text{(iii bis)} \hand{$\Longrightarrow$} \text{(iii)} \hand{$\Longrightarrow$} \text{(ii)} \]
compte tenu que pour la forme standart, le discriminant modifié est 1.
Nous prouverons maintenant (ii)\hand{$\Longrightarrow$}(iii bis), puis (i)\hand{$\Longleftrightarrow$}(ii).

Soit $\underline{E}_o = \underline{O}_S^{\ n}$ , $Q_o$ la forme standart sur $\underline{E}_o$ , Considérons le
foncteur
\[ \underline{\mathrm{Isom}}( (\underline{E}_o,Q_o),(\underline{E},Q)) :\ (\mathrm{Sch})_{/S}^{\ o} \longrightarrow (\mathrm{Ens}) , \]
dont la valeur en tout S' sur S est formé de l'ensemble des isomorphismes
$\underline{E}_{o\,S'} \to \underline{E}_{S'}$ compatibles avec les formes $Q_{o\,S'}$ et $Q_{S'}$ . Il est immédiat
(sans condition sur \inserted{$\underline{E}$ \unsure{ni}} Q) que ce foncteur est représentable par un préschéma
P affine et de présentation finie sur S, qui est un \inserted[typed]{sous-}\deleted{pré}schéma de W =
$W(\underline{\mathrm{Hom}}_{\underline{O}_S}(\underline{E}_o,\underline{E}))$ , et un sous-préschéma fermé de \deleted{$\underline{\mathrm{Isom}}(\underline{E}$} l'ouvert
$\underline{\mathrm{Isom}}\ (\underline{E}_o,\underline{E}_1)$ de W \unsure{/} L'implication (ii)\hand{$\Longrightarrow$}(iii bis) sera prouvée si
nous montrons que lorsque $\overset{\hand{$\sim$}}{d}(Q)$ est inversible, P est fidélement plat
sur S: on prendra en effet S'=P.

Posons pour simplifier \xout{\unsure{$\underline{G}(\underline{E}_o$}x}

\mspage{\hand{178}}
